% CLAUDE REPORT::WCLB-0001::2026-11-08
% !TEX program = xelatex
\documentclass[11pt]{article}
\usepackage[margin=1in]{geometry}
\usepackage{amsmath,amssymb,amsthm}
\usepackage{xurl}
\newtheorem{proposition}{Proposition}
\newtheorem{corollary}[proposition]{Corollary}
\theoremstyle{remark}
\newtheorem*{remark}{Remark}
\newcommand{\status}[1]{\textbf{[#1]}}
\newcommand{\Wb}{\overline W}
\title{CLAUDE WCLB-0001\\[2pt]\large What the bridge $\mathrm{SAT}\in\mathrm P\Rightarrow\Wb_R\le\mathrm{poly}$ can and cannot mean:\\ vacuity, the optimality reading, and a filter for candidate systems $R$}
\author{CLAUDE (scheduled worst-case run)}
\date{2026-11-08}
\begin{document}
\maketitle
\noindent\textbf{Status.} P versus NP remains OPEN. B1, B2, B3 (the bridge) remain OPEN for every concrete $R$ in the corpus.
Additive report; no schedules, ownership, or root \texttt{main.tex} changed. Nothing here is a lower bound.

\section{Corpus reviewed (scoped)}
Read: \texttt{AGENTS.md}; root \texttt{main.tex} B-ledger header; five-role handoffs LSPV-1/337 (finite reflection), LSPV-3/492 (one-hot mixed shortcut), LSPV-5/428 (distinct-support affine); \texttt{research/agents/claude/world-ledger} README. Not re-read byte by byte: older MAIN snapshots. \textbf{Duplication audit:} Prop.~3 below restates what 337 already contains (Cook--Krajicek style sufficient condition); Props.~1,2 and Cor.~1 were not found in the corpus by keyword search (``vacuous'', ``optimal proof system''), but their proofs are short and may be folklore; no priority is claimed.

\section{Definitions}
A Cook--Reckhow system $R=(V,\mathrm{cost})$: $V$ polynomial-time, $\mathrm{cost}_R(\pi)\ge|\pi|$ and polynomially bounded in $|\pi|$ (so cost is polynomially related to proof length). Write $s_R(F)=\min_{\pi:F\vdash_R\bot}\mathrm{cost}_R(\pi)$ and $\Wb_R(N)=\max_{|F|\le N,\,F\in\mathrm{UNSAT}} s_R(F)$.
\begin{proposition}[free optimal sequence is the standard measure]\status{PROVED/EXACT}
$\Wb_R$ is exactly the classical worst-case proof-size function. Not charging the search for $\pi$ is the definition of $s_R$, so the minimum over derivations is a legitimate lower-bound target; and $\Wb_R\le\mathrm{poly}$ implies $\mathrm{UNSAT}\in\mathrm{NP}$ (guess $\pi$, run $V$).
\end{proposition}

\section{Main results}
Let $\mathrm{Bridge}(R)$ denote the implication $\mathrm{SAT}\in\mathrm P\Rightarrow\Wb_R\le\mathrm{poly}$.
\begin{proposition}[vacuity]\status{PROVED/EXACT}
(a) If $\mathrm P\ne\mathrm{NP}$ then $\mathrm{Bridge}(R)$ is true for every $R$.
(b) If $\Wb_R$ is unconditionally known to be superpolynomial (e.g.\ Resolution, Cutting Planes, bounded-depth Frege, polynomial calculus over fixed fields), then $\mathrm{Bridge}(R)\iff \mathrm P\ne\mathrm{NP}$.
\end{proposition}
\emph{Proof.} (a) false antecedent. (b) If $\mathrm P\ne\mathrm{NP}$, (a). If $\mathrm P=\mathrm{NP}$, the antecedent holds and the consequent contradicts the known lower bound, so $\mathrm{Bridge}(R)$ fails. $\square$

\begin{corollary}[admissibility filter]\status{PROVED/EXACT}
A proof of $\mathrm{Bridge}(R)$ for a system with a known unconditional superpolynomial bound is a proof of $\mathrm P\neq\mathrm{NP}$. Hence a candidate $R$ is useful for the program only if $\Wb_R$ is \emph{not} already known superpolynomial; and for such $R$, $\mathrm{Bridge}(R)$ must be proved by an upper-bound (simulation) argument valid in the world $\mathrm P=\mathrm{NP}$, not by lower-bound techniques. Any lower-bound statement for $R$ plus $\mathrm{Bridge}(R)$ jointly yields $\mathrm P\ne\mathrm{NP}$ and no weaker conclusion.
\end{corollary}

\begin{proposition}[optimality reading]\status{PROVED/EXACT}
For a Cook--Reckhow $R$: $\mathrm{Bridge}(R)$ holds iff in the world $\mathrm P=\mathrm{NP}$, $R$ is polynomially bounded, i.e.\ $R$ is optimal (simulates every Cook--Reckhow system) in that world. In particular every optimal proof system satisfies $\mathrm{Bridge}$.
\end{proposition}
\emph{Proof.} If $\mathrm P=\mathrm{NP}$ with decider $A$, the system $T_A$ (empty proof, $V$ runs $A$) is polynomially bounded, so a system is polynomially bounded iff it simulates $T_A$; and any polynomially bounded system simulates every system (since $s_{R'}(F)\ge|F|$ up to padding). Conversely a simulating $R$ is polynomially bounded. $\square$
\begin{remark}
Whether an explicit optimal system exists unconditionally is a known open problem (it exists if $\mathrm{NE}=\mathrm{coNE}$, Krajicek--Pudlak 1989). So $\mathrm{Bridge}(R)$ for explicit $R$ is weaker than optimality, but not known for any explicit $R$ of Frege strength. This is the sharp form of the corpus's B3 difficulty.
\end{remark}

\begin{proposition}[monotonicity; the tension]\status{PROVED/EXACT}
If $Q$ p-simulates $R$ then $\mathrm{Bridge}(R)\Rightarrow\mathrm{Bridge}(Q)$ and $\Wb_Q\ \text{superpoly}\Rightarrow\Wb_R\ \text{superpoly}$. So the bridge is easier for strong systems and the lower bound is easier for weak ones; the target is a system at the boundary where both are provable, which by Cor.~1 can only lie among systems whose lower bounds are open.
\end{proposition}

\begin{proposition}[known sufficient route; restates LSPV-1/337]\status{KNOWN THEOREM / CONDITIONAL}
Extended Frege simulates every system whose soundness is provable in PV (Cook 1975; Krajicek--Pudlak). Hence $\mathrm{Bridge}(eF)$ would follow from: $\mathrm P=\mathrm{NP}\Rightarrow$ some polynomial SAT decider has a PV-provable (or poly-size-finite-instance) soundness proof. That premise is \status{OPEN}. Levin universal search does not supply it: it is correct on satisfiable inputs only, since unsatisfiability has no efficiently checkable witness absent the unknown decider.
\end{proposition}

\section{Literature clues (searched this run)}
\begin{itemize}
\item Bhattacharya--Chattopadhyay, arXiv:2507.23008 (STOC 2026): $\exp(\tilde\Omega(N^\epsilon))$ for $\mathrm{Res}(\oplus)$ refutations of lifted expander Tseitin formulas of depth $N^{2-\epsilon}$. Efremenko--Itsykson (ECCC TR25-188): $\exp(N^{\Omega(1)})$ for $\mathrm{Res}(\oplus)$ of depth $O(N\log N)$. General dag-like $\mathrm{Res}(\oplus)$ superpolynomial bounds remain open. \status{ANALOGY/CANDIDATE}: the project's mixed XOR/AND/RES calculus contains $\mathrm{Res}(\oplus)$-style steps, so any hard family for the full mixed calculus must in particular beat unrestricted $\mathrm{Res}(\oplus)$, which no known technique does; this explains why plain Tseitin families (LSPV-3/492, linear mixed proofs) cannot work and why lifted/IP-gadget families are the nearest literature-supported candidates. Unrestricted-depth bound would be new mathematics.
\item Cook--Krajicek and follow-ups on NP$\subseteq$P/poly provability and proof systems with advice: advice-augmented Frege systems are optimal among advice systems, which trivialises Bridge at the price of non-uniform lower bounds. \status{ANALOGY/CANDIDATE}; details not verified from primary text this run.
\end{itemize}
Classification: Props 1--4 are exact transfers within the stated model; the Res$(\oplus)$ link is an analogy needing a simulation theorem of the mixed calculus by $\mathrm{Res}(\oplus)$ or vice versa.

\section{Task consequences (no change to allocations)}
\begin{itemize}
\item \textbf{Lead / Builder:} before investing in a lower bound for $R$, record whether $\Wb_R$ is already known superpolynomial (then Cor.~1 applies and the bridge is P$\ne$NP) or open. State which Proposition~2--3 route proves Bridge.
\item \textbf{Certificates:} test families against unrestricted $\mathrm{Res}(\oplus)$ upper bounds first (cheap filter).
\item \textbf{Resource auditor / Red team:} reject any claimed Bridge proof that does not use $\mathrm P=\mathrm{NP}$ non-black-box (relativization audit pending: not done this run).
\end{itemize}
\noindent\textbf{Honest summary.} No lower bound and no bridge proof were obtained. The contribution is a precise filter: the bridge is vacuous in a P$\neq$NP world, equivalent to P$\neq$NP for any system with a known lower bound, and equals ``optimality in the P$=$NP world'' otherwise. P versus NP remains OPEN.

% END_OF_VERIFIED_REPORT::CLAUDE::WCLB-0001
\end{document}
